\section{Zadání}
Pro zvolenou dvojici států provnejte hodnoty délkového zkreslení u následujících kartografických zobrazení:
\begin{itemize}
    \item Konformní válcové zobrazení se dvěma nezkreslenými rovnoběžkami (Mercartovo zobrazení)
    \begin{align*}
x &= R \cdot \cos(u_0)\cos(v) \\
y &= R \cdot \ln\tan(\frac{u}{2}+45^\circ).
    \end{align*}
\item Konformní kuželové zobrazení se dvěma nezkreslenými rovnoběžkami (Lambertovo zobrazení)
    \begin{align*}
\rho &= \rho_0 \cdot \left(\frac{\tan(\frac{u_o}{2}+45^\circ)}{\tan(\frac{u}{2}+45^\circ)}\right) \\
\epsilon &= n \cdot v.
    \end{align*}
\item Konformní azimutální zobrazení (Stereografická projekce)
    \begin{align*}
\rho &= 2R \cdot \tan \left(\frac{90^\circ-u}{2}\right) \\
\epsilon &=  v.
    \end{align*}
\end{itemize}

Kartografická zobrazení pro každý stát navrhněte v obecné poloze, parametry potřebné pro výpočet odvoďte z vhodného mapového pokladu. 

Pro navrhovaná území volte dvě nezkreslené rovnoběžky tak, aby ve středu i na okraji zájmového území byly stejné absolutní hodnoty délkového zkreslení: $m_r^s = m_r^j = 1 +v$, $m_r^r = 1 -v$.

V závěru proveďte diskuzi nad dosaženými hodnotami, pokuste se tyto výsledky zobecnit a rozhodnout, které zobrazení je pro každý ze států nejvhodnější.
Úhlové výpočty proveďte s přesností na minuty, výpočty zkreslení na 6 desetinných míst.

\noindent
\begin{tabularx}{\textwidth}{|X|c|}
    \hline
    \textbf{Krok} & \textbf{Hodnocení} \\
    \hline\hline
    {Srovnání kartogrfikckých zobrazení pro vybrané územní celky.} \cellcolor[HTML]{94F29E} & 20b \\
    \hline
    \textit{Exaktní výpočet dotykové rovnoběžky (válcové zobrazení).} \cellcolor[HTML]{94F29E} & \textit{+5b} \\
    \hline
    \textit{Kartografický pól s využitím analytické geometrie (válcové zobrazení).} \cellcolor[HTML]{94F29E} & \textit{+10b} \\
    \hline
    \textit{Kartografický pól s využitím analytické geometrie (kuželové zobrazení).} \cellcolor[HTML]{94F29E} & \textit{+15b} \\
    \hline
    \textit{Kartografický pól jako střed opsané kružnice s min. poloměrem (azimutální zobrazení).} \cellcolor[HTML]{94F29E} & \textit{+10b} \\
    \hline
    \textit{Ekvideformáty délkového zkreslení s vhodně zvoleným krokem.}  & \textit{+x $\cdot$5 b} \\
    \hline
    \textbf{Max celkem:} & \textbf{75b} \\
    \hline
    
\end{tabularx}

\textit{Splněno \textbf{60 b} + \textbf{10 b} LaTeX.}


